\section{Starting Point}
\subsection{State of the art and preliminary work}\label{sec:stateoftheart}

\subsubsection{Planning knowledge as the bottleneck of LLM-based planning assistance}

Factory and logistics planning is a document- and expert-driven process. Requirements, process descriptions, capacity figures and layout variants are spread over specifications, tables and meeting notes, and are combined by experienced planners according to guidelines such as VDI~5200 \parencite{vdi5200,wiendahl2015handbook}. \Acp{LLM} promise to make this knowledge usable for planning assistants, for example to answer planning questions or to generate simulation models from planning documents \parencite{ferry2025anwendung,schmitt2026llm,kraemer2026automated,kute2025generative}. To ground the models in project-specific facts, recent systems combine \ac{RAG} with \acp{KG} that are built from the documents by the \ac{LLM} itself (\ac{GraphRAG}) \parencite{edge2024local,peng2026survey}. In production, \ac{KG}- or ontology-grounded \ac{RAG} has been applied to shop-floor question answering, failure mode and effects analysis, maintenance assistance and CNC process planning \parencite{wan2025hybrid,bahr2025fmea,ludwig2025ontology,hoang2025arkness}.

How reliable such assistance is depends on the quality of the knowledge it retrieves. The evidence base on this question is thin. A systematic review by the applicants' group on \acp{LLM} as decision support for strategic manufacturing and logistics planning (protocol: \cite{reider2026slrprotocol}; results in preparation) found that none of 34 included studies used an experimental design, that 32 report results only from the developers, and that reliability was the most frequently named barrier (18 of 34). A second review of 88 studies on \ac{LLM}-based simulation-model generation, conducted for the ASMG4SimPy project (\autoref{sec:prelim}), found knowledge-graph grounding in only 7 studies, none of which tested its contribution by ablation. Reliability of the knowledge layer is therefore a central, open question for planning assistance.

\subsubsection{GraphRAG: from documents to answers}

\ac{GraphRAG} pipelines consist of four knowledge-related stages that precede generation (\autoref{fig:pipeline}): a schema or ontology defines which entity and relation types exist; an extractor turns text into triples; entity reconciliation merges duplicates and separates homonyms; and a retriever selects the subgraph or summaries passed to the \ac{LLM}. Errors arise in the first three stages, but are usually measured only at the end.

\begin{figure}[t]
\centering
\begin{tikzpicture}[node distance=0.35cm, every node/.style={font=\footnotesize\sffamily},
  st/.style={draw=none, fill=blue!25, rounded corners=2pt, minimum height=1.0cm, text width=2.5cm, align=center},
  ar/.style={-{Latex[length=2mm]}, thick, black!60}]
\node[st] (s1) {Schema /\\ ontology};
\node[st, right=of s1] (s2) {KG\\ construction};
\node[st, right=of s2] (s3) {Entity\\ reconciliation};
\node[st, right=of s3, fill=black!10] (s4) {Graph\\ retrieval};
\node[st, right=of s4, fill=black!10] (s5) {Generation\\ (answer, model)};
\draw[ar] (s1)--(s2); \draw[ar] (s2)--(s3); \draw[ar] (s3)--(s4); \draw[ar] (s4)--(s5);
\draw[decorate, decoration={brace, amplitude=5pt, mirror}, thick] ([yshift=-2pt]s1.south west) -- ([yshift=-2pt]s3.south east) node[midway, below=6pt] {errors are created here (focus of this project)};
\draw[decorate, decoration={brace, amplitude=5pt, mirror}, thick] ([yshift=-2pt]s4.south west) -- ([yshift=-2pt]s5.south east) node[midway, below=6pt] {errors become visible here};
\end{tikzpicture}
\caption{Knowledge-related stages of a \ac{GraphRAG} pipeline. The project addresses the construction stages and uses retrieval and generation as the evaluation layer.}
\label{fig:pipeline}
\end{figure}

\paragraph{Schema and ontology construction.}
\acp{LLM} are used to learn ontologies from text \parencite{babaeigiglou2023llms4ol}, to generate ontologies from \acp{CQ} \parencite{lippolis2025ontology}, to guide extraction with a fixed schema \parencite{mihindukulasooriya2023text2kgbench} and to induce schemas during extraction \parencite{zhang2024edc,bai2026autoschemakg}. AutoSchemaKG induces schemas over more than 50~million documents with 92\,\% semantic alignment to human-crafted schemas \parencite{bai2026autoschemakg}. Benchmarks show, however, that \acp{LLM} recover vocabulary and taxonomies reliably but are weakest on property-mediated relations, constraints and axioms \parencite{zhang2026ontourl,li2026cq4oe}, which are exactly the elements planning needs (capacities, precedence, validity). In manufacturing, \ac{LLM} work concentrates on populating standard models such as the \ac{AAS} \parencite{xia2024aas}. Generated schemas are evaluated against gold ontologies, not by their effect on downstream answers.

\paragraph{Knowledge-graph construction.}
Construction methods range from prompted per-chunk extraction \parencite{edge2024local,gutierrez2024hipporag} through extract-then-canonicalise pipelines \parencite{zhang2024edc,mo2025kggen} and document-level fusion \parencite{zhang2025rakg} to cost-efficient indexing \parencite{guo2025lightrag,huang2025ketrag} and temporal graphs \parencite{rasmussen2025zep}; see \parencite{zhong2024survey,pan2024unifying} for surveys. Quality is rarely measured. Where it is, metrics are unstable: on the MINE benchmark, Microsoft GraphRAG is reported to recover 48\,\% of facts in one study \parencite{mo2025kggen} and 90\,\% on a rebuilt version in another \parencite{zhang2025rakg}. Validation methods filter spurious triples \parencite{huang2024graphjudge}, but missing facts are neither detected nor repaired at construction time. Graph construction from engineering documents with tables, units and numeric constraints has been demonstrated for CNC process planning, but without a gold graph \parencite{hoang2025arkness}.

\paragraph{Entity reconciliation.}
Widely used \ac{GraphRAG} systems merge entities by exact name matching \parencite{edge2024local,guo2025lightrag} or add similarity edges \parencite{gutierrez2024hipporag,gutierrez2025rag}. \ac{LLM}-based entity matching reaches an average F1 of 88.9 on record-pair benchmarks \parencite{peeters2023entity} and has been extended to selection and clustering strategies \parencite{wang2024match,fu2025incontext}, but these benchmarks consist of structured product and bibliographic records, not of free-text entities in \ac{LLM}-built graphs. Graph denoising by aggressive merging improves \ac{GraphRAG} answers up to a point and degrades them when too many entities are merged \parencite{zheng2025less}; extraction errors also bias downstream graph analyses \parencite{cai2025understanding}. In manufacturing data integration, agent-based \ac{ER} did not yield a significant gain \parencite{bega2025mixture}. A controlled comparison of split and merge errors in terms of their downstream effect has not been published outside the applicants' preliminary work.

\paragraph{Graph retrieval and evaluation.}
Retrieval is the most active strand: community summaries \parencite{edge2024local}, dual-level retrieval \parencite{guo2025lightrag}, \ac{PPR} \parencite{gutierrez2024hipporag,gutierrez2025rag} and agentic exploration have been proposed; see \parencite{peng2026survey}. Controlled comparisons published since 2025 qualify the gains. Plain \ac{RAG} matches or outperforms \ac{GraphRAG} on simple factual questions \parencite{xiang2025when}, global community search costs two orders of magnitude more tokens per query \parencite{zhou2025indepth}, and swapping the order of answers changes \ac{LLM}-judge win rates by more than 30\,\% \parencite{zeng2025significant,han2026rag}. Robust retrieval over noisy or incomplete graphs has been proposed \parencite{ma2026dontworry}, and the graph has been identified as an attack surface \parencite{liang2026underfire}. Benchmarks are drawn from Wikipedia, novels and textbooks \parencite{xiang2025when,xiao2025graphragbench}; dataset pitfalls further limit their validity \parencite{zhang2025diagnosing}.

\subsubsection{Research gap}\label{sec:gap}

Across these strands, three observations define the gap that this project addresses.
\begin{enumerate}[label=(\roman*), leftmargin=*, itemsep=1pt]
  \item \textbf{Construction errors are created upstream but measured downstream.} Schema, extraction and reconciliation methods report end-task gains for one pipeline; the effect of a specific construction error on a planning outcome is rarely isolated.
  \item \textbf{There is no validated, gold-free measure of graph quality.} Intrinsic metrics are unstable and not validated against downstream utility, and existing validation targets spurious triples rather than omissions, although omissions are the most damaging error type (\autoref{sec:prelim}).
  \item \textbf{Engineering planning documents are not represented.} No gold-anchored benchmark exists for \ac{KG} construction from planning documents with tables, units, variants and implicit expert knowledge.
\end{enumerate}
New retrieval methods are not the focus of this project: that strand is crowded, and our preliminary work shows that retriever choice explains far less of the answer quality than graph quality.

\subsubsection{Preliminary work of the applicants}\label{sec:prelim}

The project combines the work of two groups: the group for AI in production and logistics at Otto von Guericke University Magdeburg (Lang, Reider) and the Chair of Logistics Engineering at TU Dresden (Schmidt, Ludwig). \todo{confirm official chair name}

\paragraph{Controlled error injection in GraphRAG (Magdeburg).}
In a controlled study \parencite{reider2026graph}, the gold evidence of 600 multi-hop questions from 2WikiMultihopQA \parencite{ho2020constructing} was embedded in a realistic \ac{LLM}-extracted graph, and six typed construction errors (missing entities, missing edges, split entities, merged entities, corrupted relations and spurious edges) were injected at rates of 10--50\,\%. Each error was traced through two retrievers (k-hop and \ac{PPR}), two context formats and five generators (123{,}487 answers). Graph errors explained 83--97\,\% of the variance in accuracy, the retriever 8\,\% and the generator less than 1\,\%. At a 30\,\% rate, missing entities and edges reduced accuracy from 0.89 to 0.26 and 0.40; these omissions act through retrieval and turn answers into abstentions. Corrupted relations and merged entities act partly through generation and were the only errors that made models repeat false facts. Attaching source sentences to triples recovered up to 17 accuracy points and reduced the odds of adopting false facts by about 60\,\%. Real \ac{LLM} extractors missed 21--31\,\% of gold facts. This study provides the injection, attribution and gold-anchoring methodology on which \wpref{3} and \wpref{4} build; it was conducted on a general-domain dataset, which motivates the transfer to engineering documents.

\paragraph{LLM-based simulation-model generation in factory planning (Magdeburg).}
The group developed a framework and prototype for generating simulation models in factory planning with generative AI \parencite{kute2025generative} and an LLMOps methodology for facility layout planning \parencite{ma2026llmops}; related work converts engineering papers into executable tools for \ac{LLM} agents \parencite{ma2026replicating} and generates discrete-event simulation models automatically \parencite{kute2026automatic}. In the ongoing ASMG4SimPy project (manuscript in preparation), \ac{KG}-grounded multi-agent pipelines generate SimPy models from multi-modal planning documents of industrial cases, each with an expert reference model. Element-level labelling of these reference models (Fleiss' $\kappa = 0.668$) showed that only 9--48\,\% of their content is derivable from the documents; the rest is expert decision. These cases and labels are the basis of the gold-anchored testbed in \wpref{1}. On the logistics side, the group has modelled reconfigurable supply chains with \acp{KG} and ontologies \parencite{saidi2025modeling,saidi2025semantic} and combined \acp{KG} with human-centric AI for strategic decisions \parencite{rolf2022using}.

\paragraph{Ontology-based assistance and model generation (Dresden).}
The Dresden group developed an ontology-based \ac{RAG} procedure for a voice-controlled maintenance assistant, in which \ac{OWL} graph traversal enriches the retrieved context and classical and vector-based retrieval are combined and compared across several \acp{LLM} \parencite{ludwig2025ontology}. In Speech-to-jobshop, \acp{LLM} extract entities and relations from natural-language production descriptions into an ontology from which job-shop simulation models are generated; the study benchmarks several \acp{LLM} on this extraction task \parencite{ludwig2024speech}. Further work covers voice user interfaces in manufacturing logistics \parencite{ludwig2023voice}, ontology-based capture of failure knowledge on the shop floor \parencite{scharfe2024failure} and a semantic digital twin for production scheduling \parencite{kosse2024semantic}. The Dresden group has worked on metamodels for automatic simulation-model generation since 2010 \parencite{rehm2010metamodell} and provides open-source material-flow simulation \parencite{betker2026craneflow}.

\paragraph{Complementarity.}
Magdeburg contributes the error-injection and attribution methodology and industrial planning cases with expert reference models; Dresden contributes ontology engineering, ontology-grounded retrieval and user-facing assistance systems. Neither group has yet published work on entity reconciliation in \ac{LLM}-built graphs or on gold-free quality estimation; these are the new contributions of the project.
